%% sources: 02_related_work (2.4kw) + 03_theory (6.9kw) -> target ~1kw

\section{Background and related work}\label{sec:background}

\subsection{Preliminaries}\label{sec:background_prelim}

Modern interactive web maps render client-side from vector tiles, partitioning the world into a regular grid of square tiles at discrete zoom levels following the XYZ tiling scheme.
At zoom~$z$, the world is split into $2^z \times 2^z$ tiles, each identified by column~$x$, row~$y$, and zoom~$z$.
The client requests only the tiles that intersect the current viewport, fetching new tiles as the user pans or zooms.
When a request exceeds the source's \mlprop{maxzoom}, the renderer \emph{overzooms}, upscaling the nearest available tile, so a layer's effective visibility depends on both its filter and the source's zoom bounds rather than on layer bounds alone.

A vector tile encodes map data in a compact binary form.
Each tile holds \emph{source layers} (for example roads, buildings, water), and each source layer holds \emph{features}.
A feature carries geometry (point, linestring, or polygon as integer coordinates relative to the tile extent), an optional set of key-value \emph{properties} (such as \mlval{class}=\mlval{"motorway"}), and metadata.
The dominant wire format is Mapbox Vector Tiles (MVT), a Protocol Buffers encoding with a \emph{row-oriented} layout in which each feature stores its own property tags, so accessing a single property requires parsing past all others and compression operates on an interleaved byte stream rather than homogeneous columns.
\citet{tremmelMapLibreTileNext2025b} introduced MapLibre Tiles (MLT), a columnar alternative that stores each property as an independent column with type-specific encoding, which enables per-column compression and removes the per-feature tag overhead.

The visual appearance is governed by a separate \emph{style} document.
We target the MapLibre style specification, a client-side declarative format whose JSON structure is amenable to static analysis.
A style declares \emph{sources}, which name where tile data comes from and at which zoom levels it is available, and an ordered list of \emph{layers}, the central rendering unit.
Each layer references a source and source-layer, fixes a type (\mlval{fill}, \mlval{line}, \mlval{circle}, \mlval{symbol}, and others), and carries paint and layout properties together with a boolean \emph{filter} that selects which features render.
Filters and properties may be literal constants, zoom-dependent expressions (via \mlop{interpolate} or \mlop{step}), data-driven expressions (via \mlop{match}, \mlop{case}, or \mlop{get}), or combinations.
These expression trees are small programs in a domain-specific language and can be deeply nested with redundant or unreachable branches.
Some also read MapLibre's per-feature \emph{feature state} and map-wide \emph{global state}, runtime stores mutated by the host application at render time, so they cannot be folded to constants ahead of time.

\subsection{Related work}\label{sec:background_rw}

The literature relevant to joint style and data optimization is uneven, and no prior tool composes style rewriting with data-level pruning into a single open pipeline.

\subsubsection{Tile data pruning}
Several open-source tools attack data pruning directly.
Mapbox's \texttt{vtshaver} strips layers, features, and properties a style never references, but operates only on the data side and cannot rewrite the style.
\texttt{vt-optimizer} drops invisible layers and fields in both style and data, yet limits itself to visibility-based pruning.
Tippecanoe builds tilesets with density-budget feature dropping, polygon coalescing, and line simplification, the latter drawing on the cartographic generalization tradition of Douglas-Peucker \citep{douglasALGORITHMSREDUCTIONNUMBER1973} and Visvalingam-Whyatt \citep{visvalingamLineGeneralisationRepeated1993}.
These heuristics are purely data-driven and have no knowledge of the rendering workload.
Mapbox's proprietary style-optimized vector tiles\footnote{\url{https://docs.mapbox.com/help/glossary/style-optimized-vector-tiles}, last accessed 26.5.26} is the closest commercial precedent that couples style and data, but it is closed-source, API-only, and its terms of service preclude independent evaluation.
Schema-aware generators such as Planetiler, Tilemaker, and Protomaps \citep{maronVectorTilesAre2025,wallnerOpenSourceVector2022} fix column selection at build time, upstream of where our optimizations operate.

\subsubsection{Tile encoding}
A style is a declarative specification over a tabular feature set, which casts style-driven pruning as the projection and predicate pushdown long studied by query optimizers \citep{piraheshExtensibleRuleBased1992,graefeQueryEvaluationTechniques1993,armbrustSparkSQLRelational2015}.
The columnar MLT format draws on database compression: SIMD integer codecs \citep{lemireDecodingBillionsIntegers2015}, the FSST string compressor \citep{bonczFSSTFastRandom2020}, and cascaded per-block strategy selection \citep{kuschewskiBtrBlocksEfficientColumnar2023}.
The recurring finding is that no single encoding is universally optimal and data-driven selection is necessary \citep{abadiIntegratingCompressionExecution2006,dammeComprehensiveExperimentalSurvey2019,zengEmpiricalEvaluationColumnar2023}, which motivates an instance-optimized \citep{dingSageDBInstanceOptimizedData2022,dingInstanceOptimizedDataLayouts2021} trial-based encoding competition at tile granularity.
\citet{tremmelMapLibreTileNext2025b} and the COMTiles archive format \citep{tremmelCOMTILESCASESTUDY2023} advance the encoding layer but pick strategies without consulting the style, even though pruning changes column cardinalities and value distributions and so shifts the optimal encoding.
GeoParquet shares the same per-page encodings but its row-group metadata dominates at the \qtyrange{10}{100}{\kilo\byte} tile sizes we care about.

\subsubsection{Style optimization}
Style-level work is scarce.
The MapLibre validators and the \texttt{gl-style-migrate} tool transform but do not simplify, and migration often leaves redundant \mlop{case} chains and trivially true guards.
Stamen's default-stripping utility covers one pass in isolation.
The closest academic analogue lies outside the map domain: \citet{hagueCSSMinificationConstraint2019} formalize CSS rule merging as a constrained optimization, and rule-based rewriting \citep{piraheshExtensibleRuleBased1992} and equality saturation \citep{tateEqualitySaturationNew2009,willseyEggFastExtensible2021} supply the compiler-side vocabulary.
Our expression passes apply peephole rewriting, constant folding, and fixpoint iteration over a term-rewriting system \citep{ahoCompilersPrinciplesTechniques2007,baaderTermRewritingAll1999}, while structural passes draw on selectivity estimation \citep{selingerAccessPathSelection1979} and partial evaluation \citep{jonesPartialEvaluationAutomatic1993}.

\subsubsection{Benchmarking}
Prior rendering benchmarks compare \emph{systems} rather than \emph{optimization passes}: vector against raster delivery \citep{netekPerformanceTestingVector2020}, renderer families across feature counts \citep{ballaVectorDataRendering2025}, and tile-server implementations \citep{rechsteinerPerformancevergleichOpenSourceVectorTilesServerloesungenZur2024}.
We adopt their statistical methodology, bootstrap confidence intervals, warm-up handling, and metric separation \citep{georgesStatisticallyRigorousJava2007,hoeflerScientificBenchmarkingParallel2015}, rather than extending it, and complement it with a cumulative ablation that isolates each pass within a single renderer.
